\documentclass[a4paper,11pt]{article}
\usepackage[utf8]{inputenc}
\usepackage[T1]{fontenc}
\usepackage[ngerman]{babel}
\usepackage{amsmath,amssymb}
\usepackage[margin=2.3cm]{geometry}
\usepackage{enumitem}
\usepackage{fancyhdr}
\usepackage{array}
\usepackage[table]{xcolor}
\usepackage{tikz}
\definecolor{bgdark}{HTML}{1E1E1E}
\definecolor{fgtext}{HTML}{E6E6E6}
\definecolor{gridcol}{HTML}{4A4A4A}
\pagecolor{bgdark}
\arrayrulecolor{fgtext}
\AtBeginDocument{\color{fgtext}}
\renewcommand{\headrule}{\color{fgtext}\hrule height \headrulewidth}
\newcommand{\karo}[1]{\par\noindent
  \begin{tikzpicture}
    \draw[step=5mm, gridcol, very thin] (0,0) grid (\linewidth,#1);
  \end{tikzpicture}\par}

\newcommand{\diff}[1]{\hfill\normalfont\textbf{[#1]}}
\newcommand{\aufg}[2]{\subsection*{Aufgabe #1 \diff{#2}}}
\renewcommand{\arraystretch}{1.2}

\pagestyle{fancy}
\fancyhf{}
\lhead{RADM -- Boolesche Algebra \& Schaltalgebra}
\rhead{Übungsblatt 03}
\cfoot{\thepage}

\begin{document}

\begin{center}
  {\LARGE\bfseries Übungsaufgaben Boolesche Algebra \& Schaltalgebra}\\[2pt]
  {\large Relationen, Algebra, Diskrete Mathematik -- Teil 4}\\[10pt]
\end{center}

\noindent
\begin{tabular}{@{}p{0.6\textwidth}@{}p{0.4\textwidth}@{}}
  \textbf{Name:} \hrulefill & \textbf{Datum:} \hrulefill
\end{tabular}

\vspace{4pt}
\noindent\rule{\textwidth}{0.4pt}
\vspace{4pt}

\noindent\textit{Themen: Boolesche Terme, Schalttabellen, Boolesche Gesetze,
kDN/kKN aus Wertetabellen, KV-Diagramme (auch mit don't-care), Halb-/Volladdierer,
Axiome \& Dualitätsprinzip. Negation als Apostroph $a'$.}

\vspace{6pt}

% ---------------------------------------------------------------
\aufg{1}{leicht}
Welche der folgenden Ausdrücke sind syntaktisch korrekte Boolesche Terme
(Variablen $a,b,c$)?
\[
\text{(a) } (a\vee b')\wedge c \qquad
\text{(b) } a\wedge\vee b \qquad
\text{(c) } (a\wedge b)' \qquad
\text{(d) } ()' \qquad
\text{(e) } a'b'\vee c
\]
\karo{4.5cm}

% ---------------------------------------------------------------
\aufg{2}{leicht}
Erstelle die Wertetabellen (Schalttabellen) für
\[ f(a,b)=a\vee b' \qquad\text{und}\qquad g(a,b)=(a\wedge b)'. \]
\karo{5.5cm}

% ---------------------------------------------------------------
\aufg{3}{mittel}
Vereinfache mit den Gesetzen der Booleschen Algebra:
\[ f(a,b) = ab \vee ab' \vee a'b. \]
\karo{5.5cm}

% ---------------------------------------------------------------
\aufg{4}{mittel}
Die ,,Mehrheitsfunktion`` $f(a,b,c)$ ist $1$, wenn mindestens zwei der drei
Eingänge $1$ sind:
\begin{center}
\begin{tabular}{ccc|c}
$a$&$b$&$c$&$f$\\\hline
0&0&0&0\\ 0&0&1&0\\ 0&1&0&0\\ 0&1&1&1\\
1&0&0&0\\ 1&0&1&1\\ 1&1&0&1\\ 1&1&1&1\\
\end{tabular}
\end{center}
Lies aus der Wertetabelle die \textbf{kDN} und die \textbf{kKN} ab.
\karo{5.5cm}

% ---------------------------------------------------------------
\aufg{5}{mittel}
Minimiere die Mehrheitsfunktion aus Aufgabe~4 mit einem KV-Diagramm
zur disjunktiven Minimalform (DM).
\karo{6cm}

% ---------------------------------------------------------------
\aufg{6}{mittel}
Eine Funktion $f(a,b,c,d)$ ist genau für die Belegungen $1$, bei denen
($a=1$ und $b=1$) \emph{oder} ($c=0$ und $d=0$) gilt. Bestimme mit einem
KV-Diagramm die disjunktive Minimalform (DM).
\karo{6.5cm}

% ---------------------------------------------------------------
\aufg{7}{mittel}
Gegeben eine unvollständig definierte Schaltfunktion ($d$ = don't care):
\begin{center}
\begin{tabular}{ccc|c}
$a$&$b$&$c$&$f$\\\hline
0&0&0&0\\ 0&0&1&0\\ 0&1&0&0\\ 0&1&1&0\\
1&0&0&0\\ 1&0&1&1\\ 1&1&0&$d$\\ 1&1&1&$d$\\
\end{tabular}
\end{center}
Bestimme mit dem KV-Diagramm die DM (nutze die $d$-Felder geeignet).
\karo{5.5cm}

% ---------------------------------------------------------------
\aufg{8}{mittel}
\textbf{Halbaddierer.} Stelle die Wertetabelle eines Halbaddierers auf
(Eingänge $a,b$; Ausgänge Summe $s$ und Übertrag $\ddot u$) und gib die
\emph{minimierten} Schaltfunktionen für $s$ und $\ddot u$ an.
\karo{5.5cm}

% ---------------------------------------------------------------
\aufg{9}{schwer}
\textbf{Volladdierer.} Eingänge $a,b,\ddot u_1$; Ausgänge Summe $s$ und Übertrag $\ddot u_2$.
\begin{enumerate}[label=(\alph*)]
  \item Stelle die Wertetabelle auf.
  \item Gib die kDN für $s$ an.
  \item Minimiere $\ddot u_2$ mit einem KV-Diagramm.
\end{enumerate}
\karo{6.5cm}

% ---------------------------------------------------------------
\aufg{10}{schwer}
Beweise mit den \emph{Axiomen} der Booleschen Algebra (Kommutativ-, Distributiv-,
Neutralelement-, Komplementgesetze) das $0$-$1$-Gesetz
\[ a\vee 1 = 1. \]
Welches Prinzip verbindet diese Aussage mit $a\wedge 0 = 0$?
\karo{6cm}

\end{document}
